\documentclass[11pt]{article}
\usepackage[utf8]{inputenc}
\usepackage{amsmath,amssymb}
\usepackage[margin=1in]{geometry}
\usepackage{hyperref}
\emergencystretch=3em

\title{The uniform Taylor tree for $d=2$}
\author{Claude, with Daniel Murfet}
\date{2026-09-07}

\begin{document}
\maketitle
\sloppy

\begin{abstract}
Completion of Astra~\#9's R5. For the analytic family with common envelope $(C_0,L,D,\rho)$ and
every target $T$ there is a constant $K_T=\mathrm{uniformTreeConst}$, depending only on
$(\beta,b,\rho,C_0,L,p_1,p_2,T,h,k,D)$, such that for EVERY member $c$ of the family and EVERY
$N\ge1$,
\[
\Bigl|Z_c(N)-\sum_{\alpha\in\Lambda(h,k),\ \alpha<2T}N^{-\alpha}(A_\alpha\log N+B_\alpha)\Bigr|
\le K_T\,N^{-2T}(1+\log N)
\]
(\texttt{twoD\_taylor\_tree\_uniform}). The constant is the uniform cutoff constant of unit~133 at
the cutoffs of unit~129 plus the envelope bounds of unit~134 on the discarded pole terms. Zero
\texttt{sorry}/\texttt{axiom}.
\end{abstract}

\section{The things formalised}
{\small
\begin{verbatim}
theorem pos_of_mem_poleSet (h₁ h₂ k₁ k₂ M₁ M₂ : ℕ) (hk₁ : 0 < k₁) (hk₂ : 0 < k₂) (α : ℝ)
    (hα : α ∈ poleSet h₁ h₂ k₁ k₂ M₁ M₂) : 0 < α
noncomputable def coeffEnv (b ρ C₀ L β : ℝ) (h₁ h₂ k₁ k₂ D : ℕ) (α : ℝ) : ℝ :=
  canonCEnv ρ C₀ h₁ h₂ k₁ k₂ α * envMoment β α 0 (gaussEnv β L D)
    + (canonUEnv b ρ C₀ h₁ h₂ k₁ k₂ α * envMoment β α 0 (gaussEnv β L D)
      + canonVEnv b ρ C₀ h₁ h₂ k₁ k₂ α * envMoment β α 0 (gaussEnv β L D)
      + canonCEnv ρ C₀ h₁ h₂ k₁ k₂ α * envMoment β α 1 (gaussEnv β L D))
noncomputable def uniformTreeConst (β b ρ C₀ L p₁ p₂ T : ℝ) (h₁ h₂ k₁ k₂ D : ℕ) : ℝ :=
  uniformConst β b ρ C₀ L h₁ h₂ k₁ k₂ (cutoffMg k₁ p₁ (cutoffQ p₁ p₂ T))
      (cutoffMg k₂ p₂ (cutoffQ p₁ p₂ T)) D
    + ∑ α ∈ (poleSet ... at those cutoffs).filter (fun α => ¬ α < 2 * T),
      coeffEnv b ρ C₀ L β h₁ h₂ k₁ k₂ D α
theorem twoD_taylor_tree_uniform (β b p₁ p₂ ρ C₀ L : ℝ) (h₁ h₂ k₁ k₂ D : ℕ) (hβ : 0 < β)
    (hb : 0 < b) (hbρ : b < ρ) (hk₁ : 0 < k₁) (hk₂ : 0 < k₂) (hp₁ : ((h₁ : ℝ) + 1) / k₁ = p₁)
    (hp₂ : ((h₂ : ℝ) + 1) / k₂ = p₂) (c : ℕ × ℕ → ℝ → ℝ) (hcc : ∀ ij, Continuous (c ij))
    (H : ℝ → ℝ) (hH : Continuous H) (hc : ∀ ij s, |c ij s| * ρ ^ (ij.1 + ij.2) ≤ H s)
    (hC₀ : 0 ≤ C₀) (henv : ∀ s, 0 ≤ s → H s ≤ C₀ * (1 + s) ^ D * Real.exp (β * s * L))
    (T N : ℝ) (hN : 1 ≤ N) :
    |twoDAmp β b N h₁ h₂ k₁ k₂ (anaAmp c b)
        - ∑ α ∈ polesBelowGen h₁ h₂ k₁ k₂ p₁ p₂ T,
            N ^ (-α) * (canonA β h₁ h₂ k₁ k₂ (fun i j s => c (i, j) s) α * Real.log N
              + canonB β b h₁ h₂ k₁ k₂ (anaFaceU c b) (anaFaceV c b)
                  (fun i j s => c (i, j) s) α)|
      ≤ uniformTreeConst β b ρ C₀ L p₁ p₂ T h₁ h₂ k₁ k₂ D
        * (N ^ (-(2 * T)) * (1 + Real.log N))
\end{verbatim}
}

\section{Mathematics}

$Z-\sum_{\alpha<2T}=(Z-S^{\mathrm{merged}})+\sum_{\alpha\ge2T\text{ retained}}$, where the
merged sum equals the canonical pole sum over the whole retained set (unit~129 regrouping). The
first term is bounded by the uniform cutoff constant at exponent $E\ge2T$, pushed to exponent $2T$
(unit~131's monotonicity); each discarded term is at most $(|A_\alpha|+|B_\alpha|)N^{-2T}(1+\log N)$
and $|A_\alpha|+|B_\alpha|$ is bounded by the envelope constants of unit~134 (all retained poles are
positive). Summing gives $K_T$.

\section{Paper-facing meaning}

This is the strongest form of \texttt{thm:TaylorTree} for $d=2$ obtained so far: the expansion
holds with an explicit remainder constant that is uniform over all chart amplitudes
$\eta_we^{\beta s\xi_w}$ whose coefficient arrays share a weighted bound, for every $N\ge1$ (not
only eventually). In the SLT application the noncritical coordinates $w$ range over the support of
a partition of unity, and $\xi,\eta$ analytic on a common neighbourhood give exactly such a common
bound; so the per-stratum integral $\int Z_w(N)\rho_I(w)dw$ inherits the expansion with
coefficients $\int A_\alpha(w)\rho_I$, $\int B_\alpha(w)\rho_I$ once measurability in $w$ is
established (the remaining piece of the noncritical-coordinate adapter, Astra~\#6 rank~4).

\section{Lean notes}

\begin{itemize}
\item Wrapping many intermediate objects in \texttt{set} (cutoffs, coefficient functions, the
pole set, the remainder scale $R$) keeps the long statement manageable; a final \texttt{show}
restates the goal with \texttt{polesBelowGen} unfolded to \texttt{S.filter}.
\item \texttt{bound\_mono\_exp} produces \texttt{K * N \^{} (-E) * (1 + log N)}; regroup with
\texttt{← mul\_assoc} against the \texttt{set} variable \texttt{R}.
\item The unit compiled on the first attempt.
\end{itemize}

\end{document}
